% ============================================================
% CHAPTER 3: TOOLS & TECHNOLOGIES
% MCDA 5585 / 5586 Major Project Report
% Bhavik Kantilal Bhagat | A00494758 | Saint Mary's University
% ============================================================

\chapter{Tools \& Technologies}
\label{chap:tools}

This chapter surveys the tools and technologies that underpinned my internship
deliverables. The technology stack can be divided into five layers: the AWS cloud
infrastructure on which all workloads run; the observability and monitoring platform
I built on top of that infrastructure; the development toolchain I used to author and test
automation code; the AI-native tools that reshaped how I specified and
scaffolded code; and the RPA platforms whose operational health was the ultimate subject of
the monitoring work. Each layer is described at the depth necessary to contextualize the
engineering decisions documented in subsequent chapters.

% ─────────────────────────────────────────────────────────────
\section{Cloud Infrastructure (AWS)}
\label{sec:tools_aws}
% ─────────────────────────────────────────────────────────────

The IA-SRE team operates exclusively on \textbf{Amazon Web Services (AWS)}.
All production automation workloads, monitoring infrastructure, and credential
management run within a set of AWS accounts organized around Citco's
\textit{budget code} model --- six logical platforms (\texttt{rpacfs1},
\texttt{automationhub}, \texttt{bpm}, \texttt{citcoworks}, \texttt{ia-bpo}, \texttt{meridian}, \texttt{rpacfs1})
that map to distinct automation platforms and business functions. The following AWS services were central to my internship work.

\subsection{Amazon CloudWatch}

\textbf{Amazon CloudWatch} is Citco's primary observability service, ingesting
metrics, logs, and events from virtually every AWS resource in the IA environment.
During my internship, CloudWatch was used in four distinct modes.

\textit{Metrics and Alarms.} CloudWatch Metrics surfaces the time-series data emitted
by Lambda functions (invocation count, duration, error rate, concurrent executions,
memory utilization), ECS tasks (CPU and memory reservation and utilization), SQS queues
(depth, age of oldest message, number of messages sent/received), and MSK clusters
(consumer lag per topic partition). The \textbf{CloudWatch Alarm} subsystem monitors
these time series against configurable thresholds and routes breach notifications
through Amazon SNS. The IISS-487 delivery included dynamic CloudWatch alarms provisioned
via parameterized CloudFormation templates.

\textit{Logs Insights.} \textbf{CloudWatch Logs Insights} provides a SQL-like query
language for interactive log analysis across any CloudWatch Log Group. The query
language supports filtering, aggregation, field extraction with the \texttt{parse}
command, and time-bucketed statistical summaries via \texttt{stats}. I used this
extensively for Lambda cold-start analysis, error-code extraction from structured
JSON logs, and building the per-function invocation breakdown panels in the Grafana
dashboard. A representative Logs Insights query for Lambda error rate analysis is:

\begin{lstlisting}[language=SQL, caption={CloudWatch Logs Insights query for Lambda error rate by function}]
fields @timestamp, @message
| filter @message like /ERROR/
| stats count() as errorCount by bin(5m)
\end{lstlisting}

\textit{Metric Insights.} \textbf{CloudWatch Metric Insights} extends the SQL model
to metric time series, enabling \texttt{SELECT} and \texttt{GROUP BY} queries that
aggregate across all resources of a given service type. For example, a single Metric Insights
query can return the average Lambda duration grouped by function name across all
lambda functions in the account --- eliminating the need to define one panel per function.
A key limitation I discovered during dashboard development is that Metric Insights
\texttt{GROUP BY} mode supports only four aggregation functions --- \texttt{AVG},
\texttt{SUM}, \texttt{MIN}, and \texttt{MAX} --- and does not support percentile
statistics (p50, p99). Latency percentile panels were consequently implemented using
CloudWatch's standard metrics mode with a \texttt{FunctionName: ["*"]} wildcard
dimension, which returns percentile statistics for all functions that emitted data
in the selected time range.

\textit{Dashboards.} The CloudWatch Console dashboard was used as the authoritative
alarm state viewer and for initial exploratory analysis, with Amazon Managed Grafana
serving as the presentation layer for stakeholder-facing dashboards.

\subsection{AWS Lambda}

\textbf{AWS Lambda} is the primary serverless compute service for IA automation
workloads. A Lambda function inventory I conducted in April~2026 identified
\textbf{approximately 119 Lambda functions} distributed across the five budget code
environments (see Table~\ref{tab:budget_codes} in Chapter~\ref{chap:project}).
Lambda was used in my internship work in two distinct roles.

\textit{Runtime workloads.} Lambda functions implement business automation logic across
all five budget codes --- event-driven processing, API integrations, report generation
triggers, and scheduled batch jobs. These functions were the primary monitoring targets
for IISS-487.

\textit{Infrastructure provisioner.} I authored a purpose-built Lambda function as
part of the CloudFormation alarm provisioning pipeline (IISS-487). This provisioner
function dynamically discovered all Lambda functions in each budget code environment
at deploy-time and synthesized per-function CloudWatch alarm resources, enabling a
single \texttt{sam deploy} invocation to provision alarms across the entire Lambda
fleet without manual enumeration.

\subsection{Amazon ECS and EKS}

\textbf{Amazon Elastic Container Service (ECS)} hosts the containerized services that
constitute the Meridian workflow platform and other IA workloads. The total footprint
at internship time comprised \textbf{22 ECS clusters} across the five budget codes.
ECS task definitions specify the CPU and memory allocation for each container, and
ECS service auto-scaling policies govern horizontal scaling in response to load.
Container Insights --- Amazon's native telemetry agent for ECS --- was enabled on
five of the six Meridian-specific clusters on May~11--12, 2026 (IISS-487), unlocking
task-level CPU, memory, network I/O, and storage utilization metrics that were
previously unavailable to our SRE team.

\textbf{Amazon Elastic Kubernetes Service (EKS)} is used by a subset of Meridian
services and by the \texttt{automationhub} environment (one EKS cluster). EKS workloads
share the same CloudWatch Container Insights telemetry model as ECS but are addressed
through the Kubernetes node and pod metric dimensions rather than ECS task and service
dimensions.

\subsection{Amazon SQS and MSK}

\textbf{Amazon Simple Queue Service (SQS)} provides decoupled message queuing for
asynchronous workloads. In the IA platform, SQS serves as the trigger mechanism for
UiPath Orchestrator jobs (via SQS-triggered Lambda), as a buffer for BPM completion
events, and as the queue for several batch automation inputs. The IISS-487 dashboard
includes a dedicated SQS monitoring row tracking queue depth (\texttt{ApproximateNumberOfMessagesVisible})
and message age (\texttt{ApproximateAgeOfOldestMessage}) --- the latter being the
primary leading indicator for queue backlog incidents.

\textbf{Amazon Managed Streaming for Apache Kafka (MSK)} is the message backbone of
the Meridian event architecture, carrying events between the Event Routing Service and
its downstream consumers (Action Processors, the OpenSearch consumer). MSK consumer
lag monitoring --- specifically, the \texttt{kafka.consumer\_lag} metric at the
consumer group and partition level --- became a primary SRE focus following the
April~17, 2026 incident, in which 9~million messages queued on the Kafka topics over
a 7.5-hour period with no alert surfacing to the SRE team. After I deployed the dashboard,
MSK consumer lag detection improved from a 2--3 minute manual observation window to
under 30~seconds via Grafana panel alerting.

\subsection{Amazon S3 and CloudFormation}

\textbf{Amazon Simple Storage Service (S3)} provides object storage for CloudFormation
template staging, Lambda deployment packages, and automation output artifacts.
A noteworthy operational constraint I encountered during IISS-487: the Lambda function
implementing the CloudFormation alarm provisioner grew to exceed the
\textbf{64\,KB AWS Secrets Manager payload limit} as the template corpus expanded.
The resolution was to migrate the CloudFormation template payload from Secrets Manager
to an S3 object and pass an S3 URI to the Lambda provisioner --- a pattern that also
enables versioning and lifecycle management of template artifacts.

\textbf{AWS CloudFormation} is the Infrastructure-as-Code framework used for all alarm
and Lambda resource provisioning. The IISS-487 delivery used parameterized YAML
CloudFormation templates --- one template per resource type (Lambda alarms, ECS alarms,
SQS alarms, MSK alarms) --- with environment-specific parameter files that allow the
same template to be reused across all five budget codes. Four CloudFormation stacks were
deployed on May~11, 2026, establishing the alarm baseline for the monitoring platform.
AWS Serverless Application Model (SAM) was used as the local development and deployment
toolchain on top of CloudFormation.

\subsection{AWS Secrets Manager}

\textbf{AWS Secrets Manager} manages credentials for all IA platform services:
database connection strings, API keys, UiPath Orchestrator service account tokens, and
RPA process credentials. Secrets Manager supports automatic credential rotation, and
it was a rotation event that triggered the June~8, 2026 CAIS Deadline production
incident (IISS-706): a credential rotation returned an empty string value for a UCM
service account, causing the \texttt{ia-bpo} automation process to fail with an
authentication error. The resolution required patching the credential validation logic
to detect and raise an explicit error on empty-string credentials before attempting
authentication. The 64\,KB size limit of individual Secrets Manager values, encountered
independently during the CloudFormation provisioner work, is a platform constraint
that infrastructure authors must account for when storing large structured payloads.

% ─────────────────────────────────────────────────────────────
\section{Observability \& Monitoring}
\label{sec:tools_observability}
% ─────────────────────────────────────────────────────────────

The observability platform I assembled during this internship combines four AWS-native
and Amazon-managed tools into a coherent monitoring stack. Together they address the
three pillars of observability (metrics, logs, traces) plus container-level telemetry
for the ECS and EKS workloads.

\subsection{Amazon Managed Grafana (AMG)}

\textbf{Amazon Managed Grafana (AMG)} is Citco's presentation layer for
SRE dashboards. AMG is a fully managed Grafana service --- AWS provisions, patches,
and scales the Grafana backend while the team manages workspaces, data sources, and
dashboard content. The IA-SRE workspace uses \textbf{CloudWatch as its primary data
source}, authenticated via an IAM role that grants the Grafana service account read
access to metrics and log groups across all IA-platform AWS accounts.

A notable operational event during my development was the \textbf{in-flight workspace
upgrade from Grafana version~10.4 to version~12.4} on May~13, 2026 --- while the
IISS-487 dashboard was under my active construction. The upgrade introduced breaking
changes to the CloudWatch data source query model (notably around the Metric Insights
SQL editor and variable interpolation syntax) that required targeted panel revisions.
This event surfaced an important operational lesson: AMG workspace upgrades are
managed-service events that may land without advance warning and must be accounted for
in production dashboard governance.

The final IISS-487 dashboard I delivered comprised \textbf{14 rows and 110+ panels}, organized
by resource type (SQS, Lambda, ECS, MSK, EC2, Meridian-specific services) and
environment (budget code). Dashboard variables --- which I implemented as Grafana template
variables backed by CloudWatch dimension-value queries --- allow operators to filter
the entire dashboard to a specific budget code, ECS cluster, or Lambda function
without modifying the underlying panel queries.

\subsection{CloudWatch Container Insights}

\textbf{CloudWatch Container Insights} is a CloudWatch feature that installs a
managed telemetry agent (the CloudWatch agent) on ECS and EKS nodes and collects
container-level CPU, memory, network I/O, and storage metrics at both the task and
service levels. Without Container Insights, ECS metrics are available only at the
service and cluster level via the standard \texttt{AWS/ECS} namespace, which does not
expose individual task resource utilization. Container Insights publishes metrics into
the \texttt{ECS/ContainerInsights} namespace, enabling per-task and per-container
drill-down that is essential for diagnosing resource saturation in multi-service
clusters.

I enabled Container Insights on \textbf{five of the six Meridian ECS clusters}
on May~11--12, 2026 (IISS-487), providing CPU, memory, network, and storage
utilization telemetry at the task level for the first time. The sixth cluster
(\texttt{meridian-compute-cluster}) required a separate change request and was
deferred to a follow-on ticket.

\subsection{CloudWatch Logs Insights}

\textbf{CloudWatch Logs Insights} was used as my primary log query tool for
exploratory analysis, Lambda error triage, and Event Store log parsing. The Logs
Insights query language supports structured field extraction from JSON log lines using
the \texttt{parse} command, making it well-suited for querying the IA Event Store
log format, which embeds structured fields (\texttt{event\_code}, \texttt{client\_name},
\texttt{fund\_name}, \texttt{run\_id}) as top-level JSON keys. Logs Insights queries
were embedded in Grafana as panel data sources for log-derived panels, including the
per-function Lambda error count panels in my IISS-487 dashboard.

A practical limitation of Logs Insights in the Grafana context is that log-based
panels incur higher query latency than metric-based panels and are subject to
CloudWatch API rate limits at the log group level. For high-cardinality log groups
(e.g., those receiving thousands of Lambda invocations per minute), metric filters
--- which pre-aggregate log events into CloudWatch Metrics --- were used in preference
to real-time Logs Insights queries to maintain sub-second dashboard refresh performance.

\subsection{AWS X-Ray}

\textbf{AWS X-Ray} is the distributed tracing service used to capture end-to-end
request traces across Lambda functions, API Gateway, and downstream AWS service calls.
X-Ray instruments each invocation with a trace ID and records a segment for each
service boundary traversed, enabling the construction of a service map that makes
cross-service latency and error propagation visible. During the internship, X-Ray was
\textit{researched and partially evaluated} by me --- I documented the IAM role configuration and SDK
instrumentation approach --- but full production deployment across the
IA Lambda fleet was deferred to a follow-on roadmap item (IISS-489). My investigation
confirmed that X-Ray active tracing introduces a modest per-invocation overhead
(approximately 5--10~ms) and that the sampling rate must be configured carefully for
high-frequency Lambda functions to avoid unacceptable cost escalation.

% ─────────────────────────────────────────────────────────────
\section{Development Tools}
\label{sec:tools_dev}
% ─────────────────────────────────────────────────────────────

\subsection{Python, boto3, pytest, and Flask}

\textbf{Python~3} was my primary programming language for all automation and
infrastructure code produced during the internship. I chose Python because it is
the lingua franca of AWS automation via the \texttt{boto3} SDK, it is the dominant
language for IA Event Store logging integrations, and it has mature libraries for
REST API development, test-driven development, and data processing.

\textbf{boto3} is the AWS SDK for Python, providing client and resource abstractions
for every AWS service used in my internship. In the resource inventory API (IISS-507),
I used \texttt{boto3} clients for Lambda (\texttt{list\_functions},
\texttt{list\_tags}), ECS (\texttt{list\_clusters}, \texttt{describe\_clusters},
\texttt{list\_services}), EC2 (\texttt{describe\_instances}), and the Resource
Groups Tagging API (\texttt{get\_resources}) for tag-based cross-service discovery.
Tag-based discovery was required because the \texttt{CITCOWORKS} and
\texttt{MERIDIAN} environments share the ``Shared Cloud Artifacts'' BudgetCode tag,
making BudgetCode alone insufficient as a unique filter. I used the \texttt{AppName} tag
as the disambiguating fallback.

\textbf{pytest} was my testing framework for the inventory API. The test suite
reached \textbf{562 automated tests} achieving \textbf{99\% code coverage}, covering
unit tests for every service-layer function, integration tests for API route handlers,
property-based tests using \textbf{Hypothesis} for the tag-parsing and response
serialization logic, and mock-based tests for all \texttt{boto3} client interactions.
The Hypothesis library generates random valid inputs from user-defined strategies,
enabling systematic exploration of edge cases in tag-parsing functions without manual
enumeration.

\textbf{Flask} with async route handlers (\texttt{flask[async]}) was used to build
the \texttt{ia-sre-resources-inventory} REST API (IISS-507). I selected Flask over
FastAPI because the team's existing Lambda function tooling uses the
\texttt{serverless-wsgi} adapter for Lambda-backed WSGI APIs, and Flask integrates
cleanly with this adapter. I used async route handlers for the inventory endpoints
because they aggregate results from multiple \texttt{boto3} calls concurrently,
reducing end-to-end API latency by parallelizing AWS API calls within a single
request.

\subsection{AWS SAM and CloudFormation Toolchain}

The AWS Serverless Application Model (\textbf{SAM}) CLI was the primary IaC build
and deployment tool. SAM extends CloudFormation with Lambda-specific resource types
(\texttt{AWS::Serverless::Function}, \texttt{AWS::Serverless::Api}) and provides a
\texttt{sam build} / \texttt{sam deploy} workflow for packaging Lambda deployment
archives and synthesizing the underlying CloudFormation change sets. The alarm
provisioner I developed for IISS-487 used parameterized SAM/CloudFormation templates
with environment-specific parameter files (\texttt{samconfig.toml}) enabling
repeatable deployments across all six platforms with a single
\texttt{sam deploy --config-env <env>} command.

\subsection{Git and AWS CodeCommit}

\textbf{Git} is the version control system for all automation source code, with
\textbf{AWS CodeCommit} as the hosted remote repository service. CodeCommit integrates
with IAM for repository-level access control, eliminating the need for external
credentials (SSH keys or personal access tokens) by delegating authentication to the
IAM principal making the \texttt{git} command. Repositories are organized by project
code, and branch protection rules enforce pull-request-based merges to the main
branch. I developed the inventory API (IISS-507) on a feature branch using a
specification-driven workflow (described in Section~\ref{sec:tools_ai}) before
merging through a standard PR review process.

\subsection{JIRA and Confluence}

\textbf{JIRA} is Citco's project tracking tool, used for sprint planning, story
decomposition, time logging, and status management. All my work during the internship
was tracked through JIRA tickets in the \texttt{IISS} project board, with individual
stories linked to parent epics (e.g., IISS-487 for the dashboard, IISS-507 for the
inventory API). I recorded work logs at the subtask level, providing the
granular time-tracking data that underpins the 900-hour Mitacs BSI timesheet.

\textbf{Confluence} is Citco's internal wiki, used for architecture documentation,
team knowledge-base articles, RPA process runbooks, and the Lambda function inventory
Confluence page (IISS-461) produced as a deliverable alongside the automated inventory
API. Confluence pages served as the human-readable complement to the machine-queryable
inventory: the structured table format enabled non-engineering stakeholders to browse
Lambda functions by budget code and purpose without API access.

\begin{table}[ht]
\centering
\caption{Development Toolchain Summary}
\label{tab:dev_tools}
\begin{tabular}{lll}
\toprule
\textbf{Tool} & \textbf{Category} & \textbf{Primary Use} \\
\midrule
Python 3 + boto3 & Language / AWS SDK & Lambda, ECS, EC2, SQS, MSK API calls \\
Flask (async)    & Web framework      & REST API for IISS-507 inventory service \\
pytest + Hypothesis & Test framework  & 562-test suite, 99\% coverage \\
AWS SAM / CFN   & IaC toolchain      & Alarm provisioning, Lambda deployment \\
Git / CodeCommit & Version control    & Feature branches, PR-based merges \\
JIRA            & Project tracking   & Sprint planning, work-log (900 h) \\
Confluence      & Documentation      & Architecture docs, Lambda inventory page \\
\bottomrule
\end{tabular}
\end{table}

% ─────────────────────────────────────────────────────────────
\section{AI-Native Development Tools}
\label{sec:tools_ai}
% ─────────────────────────────────────────────────────────────

A distinguishing feature of this internship was my systematic integration of
AI-native development tools into the engineering workflow --- not as novelty, but as
primary productivity infrastructure. Three tools in particular shaped how I formalized
requirements, scaffolded code, and generated test suites.

\subsection{Kiro IDE}

\textbf{Kiro} is an AI-powered development environment built on the VS~Code extension
host, developed by Amazon. It introduces a \textit{spec-driven development} model in
which engineering work flows through a structured artifact lifecycle before any code
is written.

The workflow is organized around a \texttt{.kiro/specs/} directory in the project
repository. Each feature or deliverable receives a subdirectory containing three
artifacts:
\begin{itemize}
    \item \textbf{requirements.md} --- functional requirements written in a structured
          natural-language format, linked to acceptance criteria that define testability
          conditions;
    \item \textbf{design.md} --- a technical design document that translates requirements
          into component architecture, data models, API contracts, and implementation
          decisions; and
    \item \textbf{tasks.md} --- a decomposed task list derived from the design, with
          each task scoped to a single implementation unit and cross-referenced to the
          requirement it satisfies.
\end{itemize}

Kiro's \textit{agentic execution} capability allows sub-agents to pick up tasks from
the task list, implement the corresponding code, run the project's test suite, and
open pull requests --- all autonomously and without leaving the IDE context. For the
inventory API (IISS-507), this workflow reduced the time from requirements sign-off to
a green 562-test suite by enabling parallel specification and scaffolding: while I was
conducting domain analysis and requirement writing, Kiro sub-agents were already generating
stub implementations and test frameworks for the completed requirement sections.

\textit{Steering files} and \textit{hooks} are two additional Kiro mechanisms I used
during the internship. Steering files (stored in \texttt{.kiro/steering/}) inject
project-specific context --- coding conventions, naming patterns, AWS account
constraints --- into every sub-agent invocation, ensuring generated code conforms to
the team's existing standards without per-task manual instruction. Hooks are event
triggers that fire Kiro actions on filesystem events (e.g., running \texttt{pytest}
automatically when a new \texttt{*.py} file is saved), creating a continuous feedback
loop between code generation and test validation.

The net effect on my IISS-507 delivery was a substantially faster scaffolding phase
and a higher baseline test coverage at first-draft than I would have achieved
through manual TDD alone. The spec workflow also produced persistent design artifacts
(the \texttt{requirements.md} and \texttt{design.md} files) that served as the
authoritative source of truth for PR reviews and remain in the repository as
long-lived documentation.

\subsection{Amazon Q}

\textbf{Amazon~Q} (formerly Amazon CodeWhisperer) is AWS's integrated code suggestion
and chat service, available directly within Kiro and other VS~Code-based editors.
I used Amazon~Q primarily for two purposes: inline code completion for \texttt{boto3}
API calls (where it surfaces correct method signatures and parameter names without
requiring documentation lookup) and for ``memory bank'' context retrieval, in which
the assistant uses project-level context to answer questions about existing code
without me having to manually locate the relevant file. In the context of
a large multi-account AWS environment with many resource types, Amazon~Q's ability to
generate correct boto3 pagination patterns (\texttt{paginate} calls with the
appropriate \texttt{PaginationConfig}) and error handling idioms reduced routine
API-scaffolding time substantially.

\subsection{Claude Sonnet}

\textbf{Claude Sonnet~4.5} and \textbf{Claude Sonnet~4.6} (Anthropic), accessed
through the Kiro IDE interface, were used for architecture reasoning, requirement
disambiguation, and test strategy design. For tasks requiring extended chain-of-thought
--- such as designing the tag-based resource discovery fallback strategy, reasoning
through the tradeoffs between CloudWatch Metric Insights SQL mode and standard wildcard
dimension queries, or structuring the Flask async route handler pattern for concurrent
\texttt{boto3} calls --- Claude provided substantive technical discussion that
accelerated my decision-making. I also used the models to generate first-draft
CloudFormation YAML templates from prose descriptions of alarm requirements, which
I then reviewed and refined against the actual CloudWatch metric namespaces and
dimension schemas.

The combination of Kiro's structured spec workflow, Amazon~Q's inline assistance, and
Claude's extended reasoning capability effectively compressed my \textit{design-to-green-test}
cycle for IISS-507 and contributed to the quality of the CloudFormation alarm
templates I deployed in IISS-487. It is worth noting that I treated AI-generated code and
configurations as drafts requiring human review against AWS documentation
and production environment constraints --- not as production-ready artifacts. The
models' occasional hallucinations of non-existent CloudWatch metric names or incorrect
Grafana query syntax reinforced the importance of validation against authoritative
sources before deployment.

% ─────────────────────────────────────────────────────────────
\section{RPA Platforms}
\label{sec:tools_rpa}
% ─────────────────────────────────────────────────────────────

The IA-SRE team is responsible not only for building observability infrastructure but
also for the operational health of the Robotic Process Automation workloads that
execute fund administration business processes. Two RPA platforms are in production
use, each serving a distinct class of automation process.

\subsection{Blue Prism (BPM)}

\textbf{Blue Prism} --- referred to within the team as \textbf{BPM} (Business Process
Management) --- is the legacy RPA platform hosting the daily, monthly, and quarterly
fund operations automation processes. Blue Prism's execution model is based on
\textit{licensed robot nodes}: each process runs on a named robot registered to the
Blue Prism environment, with robots organized into resource pools. The production
environment (\texttt{PROD}) organizes robots into the COE (Centre of Excellence)
hierarchy under a \texttt{DataOps} folder containing Processes, Assets, Rules, and
Queues sub-objects.

The key business processes executing on Blue Prism during the internship period
included:

\begin{itemize}
    \item \textbf{VPL Pricing} (\textit{Valuation and Pricing}) --- a time-critical
          daily queue that opens at \textbf{5~PM EST} and is processed by three primary
          bots (with capacity for three additional). VPL handles fund Net Asset Value
          and pricing completion events; any processing delay after market close
          constitutes a high-severity incident.
    \item \textbf{Corporate Actions} --- three sub-workflows covering
          \textit{Mandatory} (involuntary corporate events applied to all holders),
          \textit{Voluntary} (shareholder-election events), and
          \textit{Ops Transfer} (operations transfer Corporate Actions).
    \item \textbf{CAIS Deadline} --- the Consolidated Audit Trail (CAIS) submission
          automation under the \texttt{ia-bpo} budget code, subject to a regulatory
          filing deadline. The June~8, 2026 production incident (IISS-706) was a
          CAIS Deadline failure caused by an empty credential returned from a Secrets
          Manager rotation event.
    \item \textbf{MESO} --- the Middle Office Month-End Sign-Off process
          (\texttt{ia\_meso\_process}), a Python and AWS-based process handling cash
          and position reconciliation, \AE xeo report generation, and E-Binder
          packaging for fund administration teams.
\end{itemize}

Blue Prism processes are monitored via the BPM Orchestrator UI and, for
queue-backed processes, through SQS queue depth metrics surfaced in the IISS-487
Grafana dashboard. The team is executing an ongoing migration of BPM processes to
UiPath Cloud, driven by the Blue Prism licensing cost model and UiPath's stronger
cloud-native integration story.

\subsection{UiPath Cloud and On-Premises}

\textbf{UiPath} is the target platform for new automation development and the
destination for BPM migrations. Two deployment modes are in use:

\textit{UiPath Cloud} (cloud-hosted Orchestrator) hosts the majority of active
processes. The Orchestrator folder structure mirrors the BPM model: a COE root
with \texttt{FA~DataOps} and other sub-folders containing automation processes,
assets, and queues. UiPath Cloud processes are triggered via scheduled Orchestrator
triggers, API-based job invocations, or SQS message events processed by a UiPath
Listener Lambda. The UiPath runtime's heartbeat model (periodic status updates from
the executing robot node to Orchestrator) provides the health visibility channel
monitored during support rotations.

\textit{UiPath on-premises} hosts a smaller set of processes that must run within
the Citco VPC due to network isolation requirements --- for example, processes
requiring access to internal file shares or on-premises applications not reachable
from the cloud Orchestrator. The FX Settled Status audit process is one example.

\begin{table}[ht]
\centering
\caption{RPA Platform Comparison}
\label{tab:rpa_platforms}
\begin{tabular}{lll}
\toprule
\textbf{Aspect}              & \textbf{Blue Prism (BPM)}         & \textbf{UiPath} \\
\midrule
Deployment model     & On-premises robot nodes   & Cloud Orchestrator + on-prem nodes \\
Execution model      & Licensed robot per process  & Robot units (floating) \\
Trigger types        & Scheduled, queue-based      & Scheduled, API, SQS \\
Primary use          & Legacy fund ops processes   & New builds + BPM migrations \\
Monitoring channel   & BPM Orchestrator UI + SQS   & Orchestrator heartbeat + SQS \\
Key processes        & VPL, Corp Actions, CAIS, MESO & IA-BPO processes, FA DataOps \\
\bottomrule
\end{tabular}
\end{table}

During the internship, I participated in \textbf{four RPA support rotations}
(Weeks~8, 12, 16, and 20), handling incidents across both platforms. Notable incidents
I addressed included the CAIS Deadline credential failure (June~8, IISS-706), the MESO month-end
stuck task retrigger for INNOCAP/SAMLMOS (July~1, IISS-741, classified Critical), and
multiple VPL Pricing failures related to UI selector instability and incorrect fund
bookings (IISS-753 and IISS-876). The operational experience I gained through these
rotations directly informed my alarm threshold selection and dashboard design decisions
for the IISS-487 observability platform --- each incident type corresponds to a
specific panel or alarm row in the production Grafana dashboard.
